%%%%%%%%%%%%%%%%%
% This is an sample CV template created using altacv.cls
% (v1.7.4, 30 July 2025) written by LianTze Lim (liantze@gmail.com). Compiles with pdfLaTeX, XeLaTeX and LuaLaTeX.
%
%% It may be distributed and/or modified under the
%% conditions of the LaTeX Project Public License, either version 1.3
%% of this license or (at your option) any later version.
%% The latest version of this license is in
%%    http://www.latex-project.org/lppl.txt
%% and version 1.3 or later is part of all distributions of LaTeX
%% version 2003/12/01 or later.
%%%%%%%%%%%%%%%%

%% Use the "normalphoto" option if you want a normal photo instead of cropped to a circle
% \documentclass[10pt,a4paper,withhyper,normalphoto]{altacv}

\documentclass[10pt,a4paper,withhyper]{altacv}
%% AltaCV uses the fontawesome5 and simpleicons packages.
%% See http://texdoc.net/pkg/fontawesome5 and http://texdoc.net/pkg/simpleicons for full list of symbols.

% Change the page layout if you need to
\geometry{left=1.25cm,right=1.25cm,top=1.5cm,bottom=1.5cm,columnsep=1.2cm}

% The paracol package lets you typeset columns of text in parallel
\usepackage{paracol}

% Change the font if you want to, depending on whether
% you're using pdflatex or xelatex/lualatex
% WHEN COMPILING WITH XELATEX PLEASE USE
% xelatex -shell-escape -output-driver="xdvipdfmx -z 0" sample.tex
\iftutex
  % If using xelatex or lualatex:
  \setmainfont{Roboto Slab}
  \setsansfont{Lato}
  \renewcommand{\familydefault}{\sfdefault}
\else
  % If using pdflatex:
  \usepackage[rm]{roboto}
  \usepackage[defaultsans]{lato}
  % \usepackage{sourcesanspro}
  \renewcommand{\familydefault}{\sfdefault}
\fi

% Change the colours if you want to
\definecolor{SlateGrey}{HTML}{2E2E2E}
\definecolor{LightGrey}{HTML}{666666}
\definecolor{DarkPastelRed}{HTML}{450808}
\definecolor{PastelRed}{HTML}{8F0D0D}
\definecolor{GoldenEarth}{HTML}{E7D192}
\colorlet{name}{black}
\colorlet{tagline}{PastelRed}
\colorlet{heading}{DarkPastelRed}
\colorlet{headingrule}{GoldenEarth}
\colorlet{subheading}{PastelRed}
\colorlet{accent}{PastelRed}
\colorlet{emphasis}{SlateGrey}
\colorlet{body}{LightGrey}

% Change some fonts, if necessary
\renewcommand{\namefont}{\Huge\rmfamily\bfseries}
\renewcommand{\personalinfofont}{\footnotesize}
\renewcommand{\cvsectionfont}{\LARGE\rmfamily\bfseries}
\renewcommand{\cvsubsectionfont}{\large\bfseries}


% Change the bullets for itemize and rating marker
% for \cvskill if you want to
\renewcommand{\cvItemMarker}{{\small\textbullet}}
\renewcommand{\cvRatingMarker}{\faCircle}
% ...and the markers for the date/location for \cvevent
% \renewcommand{\cvDateMarker}{\faCalendar*[regular]}
\renewcommand{\cvLocationMarker}{}


% If your CV/résumé is in a language other than English,
% then you probably want to change these so that when you
% copy-paste from the PDF or run pdftotext, the location
% and date marker icons for \cvevent will paste as correct
% translations. For example Spanish:
% \renewcommand{\locationname}{Ubicación}
% \renewcommand{\datename}{Fecha}


%% Use (and optionally edit if necessary) this .tex if you
%% want to use an author-year reference style like APA(6)
%% for your publication list
% \input{pubs-authoryear}

%% Use (and optionally edit if necessary) this .tex if you
%% want an originally numerical reference style like IEEE
%% for your publication list
% \input{pubs-num}

%% sample.bib contains your publications
% \addbibresource{sample.bib}

% Toggle visibility of Languages section (set to \showlanguagesfalse to hide)
\newif\ifshowlanguages
\showlanguagestrue

% Toggle Centric: true = show each C8 role separately (Contractor, Intern App Sec, Intern, QA); false = show one combined entry (Intern → Contractor)
\newif\ifcentricmulti
\centricmultifalse

% Toggle non-Centric experience: true = show Tech4Good and Undergraduate Course Support; false = hide them
\newif\ifshownoncentric
\shownoncentrictrue

% Toggle skill/course groups (set to \showgroupXXXfalse to hide)
\newif\ifshowgroupnetworks
\showgroupnetworkstrue
\newif\ifshowgroupsecurity
\showgroupsecuritytrue
\newif\ifshowgroupaiml
\showgroupaimltrue
\newif\ifshowgroupmisc
\showgroupmisctrue
\newif\ifshowgrouplanguages
\showgrouplanguagestrue
\newif\ifshowgroupframeworks
\showgroupframeworkstrue

% Toggle each project (set to \showprojectXXXfalse to hide)
\newif\ifshowprojecttrellix
\showprojecttrellixtrue
\newif\ifshowprojectkvm
\showprojectkvmtrue
\newif\ifshowprojectbicycle
\showprojectbicycletrue
\newif\ifshowprojectllm
\showprojectllmtrue
\newif\ifshowprojectslugquest
\showprojectslugquesttrue
\newif\ifshowprojectbloombase
\showprojectbloombasetrue

\begin{document}

% Override \cvevent to shrink location text
\renewcommand{\cvevent}[4]{%
  {\large\color{emphasis}#1\par}
  \smallskip\normalsize
  \ifstrequal{#2}{}{}{
  \textbf{\color{accent}#2}\par
  \smallskip}
  \ifstrequal{#3}{}{}{%
    {\small\makebox[0.5\linewidth][l]%
      {\BeginAccSupp{method=pdfstringdef,ActualText={\datename:}}\cvDateMarker\EndAccSupp{}%
      ~#3}%
    }}%
  \ifstrequal{#4}{}{}{%
    \par{\scriptsize\makebox[\linewidth][l]%
      {\BeginAccSupp{method=pdfstringdef,ActualText={\locationname:}}\cvLocationMarker\EndAccSupp{}%
        ~#4}%
    }}\par
  \medskip\normalsize
}

\name{Alex Asch}
\tagline{M.S. Student UCSD}
%% You can add multiple photos on the left or right
% \photoR{2.8cm}{Globe_High}
% \photoL{2.5cm}{Yacht_High,Suitcase_High}

\personalinfo{%
  % Not all of these are required!
  \email{a.asch2020@gmail.com}
  \phone{406-623-7277}
  \location{San Diego, CA}
  \linkedin{https://www.linkedin.com/in/alex-asch-755270219/}
  \github{aasch2020}
  % \orcid{0000-0000-0000-0000}
  %% You can add your own arbitrary detail with
  %% \printinfo{symbol}{detail}[optional hyperlink prefix]
  % \printinfo{\faPaw}{Hey ho!}[https://example.com/]

  %% Or you can declare your own field with
  %% \NewInfoFiled{fieldname}{symbol}[optional hyperlink prefix] and use it:
  % \NewInfoField{gitlab}{\faGitlab}[https://gitlab.com/]
  % \gitlab{your_id}
  %%
  %% For services and platforms like Mastodon where there isn't a
  %% straightforward relation between the user ID/nickname and the hyperlink,
  %% you can use \printinfo directly e.g.
  % \printinfo{\faMastodon}{@username@instace}[https://instance.url/@username]
  %% But if you absolutely want to create new dedicated info fields for
  %% such platforms, then use \NewInfoField* with a star:
  % \NewInfoField*{mastodon}{\faMastodon}
  %% then you can use \mastodon, with TWO arguments where the 2nd argument is
  %% the full hyperlink.
  % \mastodon{@username@instance}{https://instance.url/@username}
}

\makecvheader
%% Depending on your tastes, you may want to make fonts of itemize environments slightly smaller
% \AtBeginEnvironment{itemize}{\small}

%% Set the left/right column width ratio to 6:4.
\columnratio{0.6}

% Start a 2-column paracol. Both the left and right columns will automatically
% break across pages if things get too long.
\begin{paracol}{2}
\cvsection{Experience}

\ifcentricmulti
\cvevent{Software Engineering Contractor, Application Security}{Centric Software}{September 2024-June 2025}{San Jose, CA}
\begin{itemize}
\item Utilized common application security tools including Burp Suite to reproduce pentest findings.
\item Identified and repaired security vulnerabilities on the C8 platform.
\item Implemented improved Access control framework for object views and creation.
\item Re-architected login page to modern React component, and re-configured webpack for authenticated access control.
\item Remediated many email API endpoint vulnerabilities for phishing prevention and input sanitization.
\item Improved application configuration and installer to support more secure configurations.
\end{itemize}

\divider

\cvevent{Software Engineering Intern, Application Security}{Centric Software}{June 2024-September 2024, June 2025-September 2025}{}
\begin{itemize}
\item Utilized common application security tools including Burp Suite to reproduce pentest findings.
\item Identified and repaired security vulnerabilities on the C8 platform using full‑stack knowledge.
\item Triaged and responded to customer security and pentest findings.
\item Gained experience across multiple areas of the Centric PLM product.
\item Leveraged prior knowledge of Centric PLM product to respond to and remediate pentest findings.
\end{itemize}

\divider

\cvevent{Software Engineering Intern}{Centric Software}{June 2022-September 2022, June 2023-September 2023}{}
\begin{itemize}
\item Collaborated within an Agile development team, actively participating in sprint planning, daily stand-ups, and iterative delivery cycles.
\item Refactored and enhanced existing Java components to integrate with a new architectural framework, improving modularity, maintainability, and code organization.
\item Designed and implemented front-end enhancements using JavaScript to improve usability and interface functionality.
\item Contributed to end-to-end feature development, gaining hands-on experience across the full software stack from user interface to backend services.
\item Modified existing Java code to fit a new framework and increase modularity.
\item Implemented JavaScript UI changes.
\end{itemize}

\divider

\cvevent{QA Intern}{Centric Software}{June 2021 - September 2021}{San Jose, CA}
\begin{itemize}
\item Worked with Docker and QA tooling to validate web product functionality.
\item Gained experience with usage of the Centric PLM product and domain.
\end{itemize}

\else
\cvevent{Software Engineer (Intern → Contractor), Application Security}{Centric Software}{Summers 2021--2023; June 2024 -- September 2025}{San Jose, CA}
\begin{itemize}
\item Utilized common application security tools including Burp Suite to reproduce pentest findings.
\item Identified and repaired security vulnerabilities on the C8 platform.
\item Implemented improved Access control framework for object views and creation.
\item Re-architected login page to modern React component, and re-configured webpack for authenticated access control.
\item Remediated many email API endpoint vulnerabilities for phishing prevention and input sanitization.
\item Improved application configuration and installer to support more secure configurations.
\item Triaged and responded to customer security and pentest findings.
\item Gained experience across multiple areas of the Centric PLM product.
\item Leveraged prior knowledge of Centric PLM product to respond to and remediate pentest findings.
\item Collaborated within an Agile development team, actively participating in sprint planning, daily stand-ups, and iterative delivery cycles.
\item Refactored and enhanced existing Java components to integrate with a new architectural framework, improving modularity, maintainability, and code organization.
\item Designed and implemented front-end enhancements using JavaScript to improve usability and interface functionality.
\item Contributed to end-to-end feature development, gaining hands-on experience across the full software stack from user interface to backend services.
\item Modified existing Java code to fit a new framework and increase modularity.
\item Implemented JavaScript UI changes.
\item Worked with Docker and QA tooling to validate web product functionality.
\item Gained experience with usage of the Centric PLM product and domain.
\end{itemize}

\fi

\divider
\ifshownoncentric
\cvevent{Tech4Good Components Team}{Tech4Good}{March 2023 - June 2023}{Santa Cruz, CA}
\begin{itemize}
\item Gained experience working with Angular component framework.
\item Created Angular components for the Annota project.
\item Gained experience working in software development teams.
\item Completed 4 distinct UI components in the course of one quarter.
\end{itemize}
\divider
\cvevent{Tech4Good Containers Team}{Tech4Good}{June 2023 - December 2023}{Santa Cruz, CA}
\begin{itemize}
\item Gained experience working with Angular container framework.
\item Created Angular containers for the Causeway and ExploreCareers projects.
\item Gained experience working in software development teams remotely.
\end{itemize}
\divider
\cvevent{Tech4Good Annota Research Dev Team}{Tech4Good}{January 2023 - March 2023}{Santa Cruz, CA}
\begin{itemize}
\item Unpaid student research position contributing to an existing project.
\item Developed software and fixed defects on the Annota collaborative annotation platform.
\item Used Angular with a Firestore backend to contribute features.
\item Participated in project research and design discussions.
\end{itemize}
\divider
\cvevent{Undergraduate Course Support}{Tutor and Grader for Computer Science Courses}{January 2023 – March 2024}{Santa Cruz, CA}
\begin{itemize}
\item Tutored and graded coursework for Introduction to Computer Networking (CSE 150) and Introduction to Database Systems (CSE 180).
\item Provided code reviews and exam preparation for networking and database systems.
\item Worked with Mininet virtual machines to evaluate network topology code.
\item Reviewed and graded student SQL projects.
\end{itemize}
\fi
\cvsection{Projects}

\ifshowprojecttrellix
\cvevent{Trellix Software TLS Analysis}{University of California San Diego}{June 2025}{San Diego, CA}
\begin{itemize}
\item Traced what Trellix HX monitoring agents send off-host using bpftrace to capture traffic despite TLS and custom crypto; mapped endpoints and payloads; compared Trellix terms to Norton 360.
\item Designed and used bpftrace probes to trace custom libssl/libcrypto and liblzma in xAgent and Qualys; recovered unencrypted payloads and mapped endpoints (xAgent /poll, config; Qualys CAPI, fragment/finalize, Manifest).
\item Recovered and interpreted the Qualys SQLite manifest to list collected commands (ps, ls, netstat, /etc/shadow, AWS env, docker ps, etc.); clarified scope of exfiltrated data (processes, files, network, mounts, kernel params, services).
\end{itemize}
\fi
\ifshowprojectkvm
\cvevent{Enhancing KVM Page Eviction: From FIFO to Approximate LRU}{University of California San Diego}{June 2025}{San Diego, CA}
\begin{itemize}
\item Replaced KVM's FIFO shadow-page eviction with Approximate LRU; second-chance clock LRU with access recency and lightweight per-page metadata; made thread-safe for multi-vCPU.
\item Modified KVM MMU code (mmu.c, mmu\_zap\_oldest\_pages()) collaboratively; helped with benchmark setup, flamegraph analysis, and interpreting results.
\item Evaluated with stress-ng (fewer page faults) and redis-benchmark (lower latency under load); flamegraphs showed fewer samples in key page-fault paths; shadow MMU path only.
\end{itemize}
\fi
\ifshowprojectbicycle
\cvevent{Computer Vision Driven Automated Bicycle Braking}{University of California San Diego}{December 2025}{San Diego, CA}
\begin{itemize}
\item OptiBrake: front camera with Lucas-Kanade optical flow to detect approaching obstacles; closed-loop controller runs brake motor until wheel speed zero, then reverses to unspool cable.
\item Worked on wheel odometry, closed-loop brake control, and hardware integration; interrupt-based debounced wheel speed (weighted average of last 9 RPM); helped design and analyze indoor trials.
\item Showed automated brake actuation from single-camera CV on Raspberry Pi in controlled trials; documented false-positive and outdoor limits (e.g., would need more compute or LIDAR for robustness).
\end{itemize}
\fi
\ifshowprojectllm
\divider
\cvevent{LLM Based Security Framework for Chatbot Agents}{University of California San Diego}{September 2025}{San Diego, CA}
\begin{itemize}
\item AccessControl Agent: separate LLM predicts tools needed from user prompt and outputs XML policy; reference monitor enforces so agent may only call approved tools in that order.
\item Created the reference monitor and integrated it with the AI agent; wrote code that converts the LLM's XML file into the state tree; assisted in prompt engineering and XML format, debugging, and deliverables.
\item Privilege separation (LLM sees only prompt; agent restricted by monitor); monitor caught many injections (e.g., 7/8 Llama, 3/13 GPT 4o); evaluated on simple, single-tool, and complex/conditional prompts.
\end{itemize}
\fi
\ifshowprojectslugquest
\divider
\cvevent{Slugquest Project}{University of California Santa Cruz}{January 2024 - March 2024}{Santa Cruz, CA}
\begin{itemize}
\item Worked on creating a web app for collaborative and gamified To-Do lists, with calendar support.
\item Implemented Golang-based backend CRUD API using the Gin web framework, along with SQLite database server.
\end{itemize}
\fi
\ifshowprojectbloombase
\divider
\cvevent{Bloombase Project}{University of California Santa Cruz}{March 2023 - June 2023}{Santa Cruz, CA}
\begin{itemize}
\item Built a web app integrating the iNaturalist API, Py4Web backend, and Google Maps front-end.
\item Integrated third-party iNaturalist API for biodiversity data; developed Py4Web backend for data processing and storage; built Google Maps-based frontend for geospatial visualization.
\item \href{https://github.com/BloomBase183/BloomBase}{GitHub Repository Link}
\end{itemize}
\fi

\switchcolumn


\cvsection{Skills}

\ifshowgroupnetworks
\cvtag{Wireshark}
\cvtag{Mininet}
\cvtag{QEMU}
\cvtag{KVM}
\cvtag{Docker}
\cvtag{Kubernetes}
\cvtag{Git}
\cvtag{CI/CD}
\cvtag{MySQL}
\cvtag{MPI}
\cvtag{GDB}
\cvtag{Kprobes}
\cvtag{bpftrace}
\cvtag{Flamegraphs}
\divider
\fi

\ifshowgroupsecurity
\cvtag{Burp Suite}
\cvtag{Ghidra}
\divider
\fi

\ifshowgroupaiml
\cvtag{OpenCV}
\cvtag{CUDA}
\divider
\fi

\ifshowgroupmisc
\cvtag{REST APIs}
\cvtag{HTML}
\cvtag{CSS}
\cvtag{XML}
\cvtag{Firebase}
\divider
\fi

\ifshowgrouplanguages
\cvtag{Python}
\cvtag{Java}
\cvtag{C/C++}
\cvtag{Go}
\cvtag{Rust}
\cvtag{TypeScript}
\cvtag{SQL}
\cvtag{Haskell}
\divider
\fi

\ifshowgroupframeworks
\cvtag{Angular}
\cvtag{Node.js}
\cvtag{React}
\cvtag{Webpack}
\fi


\ifshowlanguages
\cvsection{Languages}

\cvskill{English}{5}
\divider

\cvskill{Japanese}{3.5}

\fi
%% Yeah I didn't spend too much time making all the
%% spacing consistent... sorry. Use \smallskip, \medskip,
%% \bigskip, \vspace etc to make adjustments.
\medskip

\cvsection{Education}
\cvevent{M.S.\ in Computer Science}{University of California San Diego}{Sept 2024 -- Ongoing}{GPA 3.975/4.0}

\ifshowgroupnetworks
\cvtag{Operating Systems}
\cvtag{Networked Services}
\cvtag{Parallel Computation}
\cvtag{Intro to Embedded Computing}
\cvtag{Virtualization and Cloud Computing}
\cvtag{Principles of Database Systems}
\cvtag{Database System Implementation}
\cvtag{Exploring Architecture}
\divider
\fi

\ifshowgroupsecurity
\cvtag{Graduate Computer Security}
\cvtag{Language-based Security}
\divider
\fi

\ifshowgroupaiml
\cvtag{Probabilistic Reasoning and Learning}
\cvtag{Systems for LLMs and AI Agents}
\cvtag{Machine Learning for Music}
\divider
\fi

\ifshowgroupmisc
\cvtag{Algorithm Design and Analysis}
\cvtag{Principles of Programming Languages}
\cvtag{The Legacy of the Cypherpunks}
\cvtag{Programmers are People Too}
\fi

\divider

\cvevent{B.S.\ in Computer Science}{University of California Santa Cruz}{Sept 2020 -- June 2024}{GPA 3.85/4.0}

\cvachievement{\faTrophy}{Graduated Cum Laude}{Awarded to top 15\% of graduating class by GPA}


% \divider




\end{paracol}


\medskip



% Adapted from @Jake's answer from http://tex.stackexchange.com/a/82729/226
% \wheelchart{outer radius}{inner radius}{
% comma-separated list of value/text width/color/detail}

% use ONLY \newpage if you want to force a page break for
% ONLY the current column
\newpage


%% Switch to the right column. This will now automatically move to the second
%% page if the content is too long.


\end{document}
